% Corrigé intégré automatiquement pour la banque Exercices.
% Source : A_Integrer/DDS_04/080_CorrecteurP/080_CorrecteurP.tex
\begin{corrigebox}[Corrigé]
\CorrigeQuestion{1}
La résolution consiste à identifier les données et l'inconnue, choisir la loi du cours adaptée, écrire l'équation symbolique avant toute application numérique, puis vérifier l'homogénéité, le signe et l'ordre de grandeur du résultat. Les hypothèses utilisées doivent être explicitement contrôlées à la fin.
\CorrigeQuestion{2}
La résolution consiste à identifier les données et l'inconnue, choisir la loi du cours adaptée, écrire l'équation symbolique avant toute application numérique, puis vérifier l'homogénéité, le signe et l'ordre de grandeur du résultat. Les hypothèses utilisées doivent être explicitement contrôlées à la fin.
\CorrigeQuestion{3}
On factorise la fonction de transfert sous forme canonique. Chaque gain, intégrateur, zéro et pôle apporte sa pente et sa phase ; les contributions sont ensuite sommées. La marge de phase se lit à la pulsation où le gain vaut $0$ dB et la marge de gain à la pulsation où la phase vaut $-180^\circ$.
\end{corrigebox}
